% ===========================================================================
% SOURCES: docs/claude_project/01_PROJECT_BRIEF.md (a snapshot of 2026-09-24:
%            its status table and its "field contact" dataset are out of date)
%          docs/design_decisions.md (DD-004, the tiered scope; DD-036, the
%            scope as amended - the dataset and the comparison)
%
% WRITE THIS LAST but one. The introduction promises what the thesis delivers,
% and you cannot write the promise honestly until chapters 5-10 exist.
% ===========================================================================
\chapter{Introduction}
\label{chap:introduction}

\section{Motivation}
\label{sec:motivation}
% Vineyard disease detection as an agricultural problem; why inference has to
% happen in the field rather than in a datacentre (connectivity, latency,
% power); why an FPGA rather than a CPU, a GPU or a phone.
% Ground this in [morellos2024iot] - the IoT grapevine-monitoring service is
% the application context this thesis accelerates.
% The CPU-or-GPU question is no longer only rhetorical: this thesis MEASURES
% it, running the same model on a Raspberry Pi 4 and a Jetson Nano - the
% boards such deployments actually use [karim2024grapeedge] - with the same
% test images and the same meter (\DD{036}). Say here only that it is
% measured; the numbers belong to Section~\ref{sec:comparison}.

\section{Objectives and Scope}
\label{sec:objectives}
% State the tiered scope from DD-004 explicitly and early - Foundation / Core /
% Stretch. It is unusual to publish one's own success criteria in a thesis and
% it is the honest thing to do here, because the tiering is what made the
% result defensible when the board work ran late.
%
% State also the deliberate constraint: hand-written Verilog-2001, not HLS and
% not a vendor IP core. Say why (layer-level control and understanding), and
% forward-reference Section~\ref{sec:hls} where the alternative is discussed.
%
% Then the two scope statements of \DD{036}, both early:
%   - the comparison with a Raspberry Pi 4 (CPU) and a Jetson Nano (GPU) is an
%     objective of the thesis, not an afterthought - the same model, images
%     and measurements on all three platforms;
%   - the data is PlantVillage's grape leaves, split by leaf (\DD{034}), and
%     nothing else: no field photographs. State it as a limit HERE, so the
%     conclusion repeats it instead of discovering it.

\section{Contributions}
\label{sec:contributions}
% Claim only what is measured. A defensible list:
%   1. an int8 MobileNetV2 accelerator, hand-written RTL, \thDspAccelAxi DSPs,
%      bit-exact against its software reference over \thElements elements;
%   2. a verification method built on golden vectors that makes "close enough"
%      an impossible answer (Chapter~\ref{chap:verification});
%   3. a measured timing-closure progression, \MHz{\thFmaxBase} to
%      \MHz{\thFmaxPhaseFour} out of context, each step attributed
%      (Chapter~\ref{chap:timing});
%   4. a complete PS integration - DMA, CSR, block design, bit-exact C
%      preprocessing - closing at \MHz{\thClkBoard} in context
%      (Chapter~\ref{chap:integration}), and running on the ZCU104;
%   5. a measured comparison of the same model, on the same test images,
%      against a Raspberry Pi 4 and a Jetson Nano: latency, power and energy
%      per inference (Section~\ref{sec:comparison}, \DD{036});
%   6. only if it runs on the board by submission: the live camera path - a
%      USB webcam to a class, with the resize bit-exact against Pillow
%      (Section~\ref{sec:camera-path}, \DD{035}).
% A board result exists: on 2026-10-07 the bare-metal application reproduced
% the known answer on the ZCU104, bit for bit (Section~\ref{sec:board-remaining}).
% Claim exactly that - one image, one run - and no more until the other
% bring-up checks are recorded.

\section{Thesis Outline}
\label{sec:outline}
% One short paragraph per chapter. Write it from the final table of contents,
% not from this skeleton.
